\section{Introduction}

\subsection{Problem Statement and Motivation}

Artificial intelligence (AI) and deep learning have made significant improvements regarding
fast and reliable decision making on a large set of information.
This ability is esppeccialy helpful when trying to detect specific patterns
The focus of this literature review is on time‑series deep learning models (1D-CNN/LSTM/TCN)
applied to wearable seizure detection/prediction using
non‑EEG wearable/autonomic biosignals: Electrocardiography (ECG), Heart Rate Variability (HRV), 
Photoplethysmography (PPG), Electrodermal Activity (EDA), and Accelerometry (ACC).
Extended temporal context, novel deep learning architectures, training strategies, 
and personalized fine-tuning can improve performance metrics
(sensitivity, AUC, f1-score, false alarms per hour (FAR/h), latency) in clinical and ambulatory contexts.

Epilepsy affects roughly 7.6 per 1,000 people and about 30\% of
patients are drug-resistant, creating demand for timely, low‑burden
monitoring and early intervention \parencite{beghiEpidemiologyEpilepsy2019,
kwanEarlyIdentificationRefractory2000,chenNewEraPersonalised2020}.
While EEG remains the reference in controlled settings, 
it is bulky, uncomfortable, and requires complex setup, and therefore not suitable for ambulatory monitoring.
Wearables offer the possibility of
continuous monitoring in ambulatory settings \parencite{chenNewEraPersonalised2020}.

Healthcare constraints directly affect modeling objectives. 
Models must achieve high sensitivity while controlling false alarms per hour (FAR/h) to avoid alarm fatigue. 
They must also balance predictive performance with interpretability and safety requirements demanded in regulated clinical contexts. 
Robustness to domain shift across patients, devices, sites, and activities is required, 
and evaluations should address data scarcity by prioritizing prospective, multimodal, patient-preserving study designs.

This review covers time-series deep learning approaches to a focused medical use case,
comparing traditional feature-based models and emerging deep temporal
architectures (1D-CNN, LSTM, TCN, attention/transformers) for seizure
detection and preictal prediction from non-EEG wearable biosignals.

The key question is which modeling and evaluation choices deliver
clinically viable, low FAR/h, high-sensitivity performance under
pseudo-prospective and prospective, ambulatory conditions.

Cross‑dataset transfer is essential to assess robustness and domain
shift \parencite{yangMultimodalAISystem2022,andradePerformanceSeizurePrediction2024}.

\subsection{Research Aim and Questions}
This literature review focuses on machine/deep learning models for seizure detection and prediction
using non-EEG wearable/autonomic biosignals, and aims to identify the optimal strategies to achieve clinically
meaningful performance in ambulatory contexts.

\vspace{0.5em}

\subsubsection{Primary Research Question}
In non‑EEG wearable seizure detection and prediction
(ECG/HRV, PPG, EDA, ACC), which modeling and evaluation choices
reported in 2015--2025 are associated with low alarm‑based false‑alarm
rates and high sensitivity with acceptable latency in ambulatory
settings?

\subsubsection{Sub‑questions}
\begin{enumerate}
  \item What ranges of sensitivity, FAR/h, and latency are reported
    under retrospective, pseudo‑prospective, and prospective designs,
    and how do alarm‑based results differ from sample‑based metrics?
  \item Across studies, how do traditional feature‑based models compare
    to 1D‑CNN/TCN/LSTM, and what evidence exists for attention/
    transformer variants on these signals, including any reported
    computational/latency considerations?
  \item How do temporal choices (window length, stride, context length)
    correlate with reported alarm‑based performance in real‑time use?
  \item What evidence is available on robustness to domain shift
    (patient/device/site/activity), including cross‑dataset/external
    validation and how do these affect alarm‑based metrics?
  \item What dataset gaps and class imbalance issues are reported, and
    what is the measured impact of augmentation or synthetic data on
    alarm‑based outcomes?
\end{enumerate}

\vspace{0.5em}
\noindent\textit{Contribution.} 
The goal is to synthesize model designs for wearable seizure detection/prediction, 
analyze how evaluation choices affect reported performance and derive practical
recommendations to achieve low FAR/h while keeping high sensitivity in
ambulatory deployment.
